\centerline{\bf \huge 11 класс}

\centerline{\normalsize{\it Работа рассчитана на \bf{240} минут}}

\problem{11.1}
Что больше, $\sin 1 \sin 2$ или $\sin 7+\sin 8$ ?


\problem{11.2}
На Новый Год магазин подарков украсил витрину гирляндой из $2 n+1$ лампочки. Все лампочки имеют попарно различные цвета, а первая и последняя лампочка не соединены напрямую. В новогоднюю ночь произошёл сбой: каждая лампочка поменяла свой цвет на цвет одной из своих соседок. Докажите, что теперь найдутся две лампочки одного цвета.


\problem{11.3}
Барон Мюнхгаузен имеет репутацию лжеца. Недавно он утверждал, что знает такое натуральное число, которое имеет не менее 2023 делителей, отличных от единицы, каждые два из которых имеют наибольший общий делитель, больший 2023. Обманывает ли нас Барон на этот раз?


\problem{11.4}
Таблица размера $n \times n$ заполнена произвольными положительными числами. Расул Владимирович подсчитал сумму чисел в каждой строке и получил числа $a_1, a_2, \ldots, a_n$. Даяна подсчитала сумму чисел в каждом столбце и получил числа $b_1, b_2, \ldots, b_n$. Затем они составили многочлены

$$
\begin{aligned}
& P(x)=\left(x-a_1\right)\left(x-a_2\right) \cdot \ldots \cdot\left(x-a_n\right) \\
& Q(x)=\left(x-b_1\right)\left(x-b_2\right) \cdot \ldots \cdot\left(x-b_n\right)
\end{aligned}
$$


В конце концов Анджур Аркадьевич решил уравнение $P(x)-Q(x)=0$ и сказал, что получил $(n-1)$ корень. Докажите, что Анджур Аркадьевич ошибся.

%1-4 Бурятия 2023 

\problem{11.5}
На горизонтальном полу лежат три волейбольных мяча радиусом 18, попарно касающиеся друг друга. Сверху положили теннисный мяч радиусом 6, касающийся всех трёх волейбольных мячей. Найдите расстояние от верхней точки теннисного мяча до пола. (Все мячи имеют форму шара.)

%11.5 Москва 2022